\documentclass[11pt,a4paper]{article}

\usepackage[margin=1in]{geometry}
\usepackage[T1]{fontenc}
\usepackage{lmodern}
\usepackage{microtype}
\usepackage{amsmath, amssymb}
\usepackage{booktabs}
\usepackage{longtable}
\usepackage{array}
\usepackage{graphicx}
\usepackage{caption}
\usepackage{subcaption}
\usepackage[table]{xcolor}
\usepackage{titlesec}
\usepackage{enumitem}
\usepackage{fancyhdr}
\usepackage{url}
\usepackage[hidelinks,bookmarks=true,bookmarksopen=true]{hyperref}
\usepackage{siunitx}
\usepackage{parskip}

\sisetup{group-separator={,}, detect-weight=true, detect-family=true}
\renewcommand{\arraystretch}{1.15}

\definecolor{pass}{RGB}{0,140,0}
\definecolor{warn}{RGB}{200,120,0}
\definecolor{fail}{RGB}{180,0,0}
\definecolor{note}{RGB}{60,60,180}
\definecolor{accent}{HTML}{01696F}
\definecolor{ink}{HTML}{28251D}
\definecolor{muted}{HTML}{7A7974}
\definecolor{rule}{HTML}{D4D1CA}

\newcommand{\PASS}[1]{\textcolor{pass}{\textbf{#1}}}
\newcommand{\WARN}[1]{\textcolor{warn}{\textbf{#1}}}
\newcommand{\FAIL}[1]{\textcolor{fail}{\textbf{#1}}}
\newcommand{\NOTE}[1]{\textcolor{note}{\textit{#1}}}

\titleformat{\section}
  {\normalfont\large\bfseries\color{accent}}{}{0pt}{}
  [\vspace{-2pt}\textcolor{rule}{\rule{\linewidth}{0.4pt}}]

\titleformat{\subsection}
  {\normalfont\normalsize\bfseries\color{ink}}{}{0pt}{}

\titlespacing*{\section}{0pt}{14pt}{6pt}
\titlespacing*{\subsection}{0pt}{10pt}{2pt}

\pagestyle{fancy}
\fancyhf{}
\renewcommand{\headrulewidth}{0pt}
\renewcommand{\footrulewidth}{0pt}
\fancyfoot[L]{\footnotesize\color{muted}use4-learning-lab --- reference run 2026-06-11}
\fancyfoot[C]{\footnotesize\color{muted}\leftmark}
\fancyfoot[R]{\footnotesize\color{muted}\thepage}

\title{\textbf{USE4 on Public Data --- Deviations Register}\\[2pt]
\large Every forced deviation, the chosen solution, and why it is defensible\\[4pt]
\normalsize USE4 Learning Lab --- Reference Audit}
\author{Audit Date: 2026-06-11}
\date{}

\begin{document}
\maketitle
\tableofcontents
\newpage

%──────────────────────────────────────────────────────────────────────────────
\section{Purpose and Scope}

This document is the centrepiece of the lab: the pipeline-level companion to the per-component
documents (\texttt{estu\_audit.tex}, \texttt{daily\_summary.tex}, the
twelve style \texttt{<factor>\_factor\_audit.tex} files, and
\texttt{industries\_factor\_audit.tex}). Those audits answer
\emph{``was each build executed correctly?''}; this one answers
\emph{``where does the pipeline deviate from Barra USE4, what did we decide,
and why is that decision defensible?''}

Every deviation from the \emph{Barra US Equity Model (USE4)} documentation
(Methodology Notes; Empirical Notes, September 2011) is registered here with a
four-part treatment: the USE4 requirement, the deviation, the chosen solution,
and the validity argument --- closing with an explicit \textbf{revisit trigger}.
The register is organised by theme: data availability (\S4), methodology
judgment calls (\S5), and engineering conventions (\S6). Deviations shared by
several factors appear once.

\smallskip
\textbf{Run audited.} All numbers are verbatim from the full pipeline run of
2026-06-11 (\path{your end-to-end runner}; data through 2026-06-10).
Every component passed its full validation battery on this run.

%──────────────────────────────────────────────────────────────────────────────
\section{Pipeline Overview}

\textbf{Dependency order:}\quad
\texttt{estu} $\to$ \texttt{daily} $\to$
\{\texttt{size}, \texttt{bp}, \texttt{dyld}, \texttt{eyld}, \texttt{gro},
\texttt{lev}, \texttt{liq}, \texttt{beta}, \texttt{mom}, \texttt{resvol}\};\quad
\texttt{beta} $\to$ \texttt{nlb};\quad \texttt{size} $\to$ \texttt{nls};\quad
\texttt{estu} $\to$ \texttt{industries}.

%──────────────────────────────────
\subsection{The reference run}

The numbers in this register come from the lab's reference run of 2026-06-12
(Sharadar data through 2026-06-10): all sixteen builds executed sequentially
in dependency order, every validation green. Your rebuild will differ in row counts
and tail statistics with data vintage; the structural signatures (stability
coefficients, monotone quintiles, standardisation identities, complete industry
coverage) should not.

%──────────────────────────────────────────────────────────────────────────────
\section{Headline Fidelity Gaps}

Six deviations are material enough to frame the whole register. All six are
\emph{data-availability} gaps --- Sharadar carries no analyst estimates, no
float shares, and no preferred-equity balance --- so each is a forced choice,
not a methodological preference, and each carries a concrete procurement-based
revisit trigger.

\begin{enumerate}[leftmargin=*,topsep=2pt,itemsep=4pt]
  \item \textbf{EPFWD omitted from EYLD (75\% of the intended composite).}
    USE4: $\mathrm{EYLD} = 0.75\,\mathrm{EPFWD} + 0.15\,\mathrm{CETOP} +
    0.10\,\mathrm{ETOP}$. With no forward-estimates feed, the dominant
    descriptor is dropped and the remainder renormalised (\S4, entry A1).
  \item \textbf{EGRLF proxied in GRO (70\% of the intended composite).}
    USE4: $0.70\,\mathrm{EGRLF} + 0.20\,\mathrm{EGRO} + 0.10\,\mathrm{SGRO}$;
    the long-term analyst growth descriptor is proxied by trailing EGRO
    (\S4, entry A2).
  \item \textbf{Total instead of float-adjusted market cap}, pervading ESTU
    selection and weights, SIZE/NLS, the benchmark, and every
    $\sqrt{\text{mcap}}$ regression weight (\S4, entry A3).
  \item \textbf{Preferred equity unavailable (PE $=0$)}, affecting BP's book
    equity and LEV's MLEV/BLEV (\S4, entry A4).
  \item \textbf{GICS unavailable.} USE4's 60 industry factors are built on a
    licensed taxonomy; the industry layer is a hand-engineered 55-factor scheme
    over Sharadar's classification atoms (\S4, entry A9).
  \item \textbf{No business-segment data.} USE4's fractional industry exposures
    ($\sim$63\% of cap weight multi-exposed) are unbuildable; every stock gets
    single 0/1 membership (\S4, entry A10).
\end{enumerate}

The common defence pattern: where USE4's \emph{inputs} are unobtainable, the
pipeline preserves USE4's \emph{construction} (PIT discipline, trim,
standardisation, orthogonalisation) on the best available proxy, and the
validation batteries confirm the resulting factors behave as USE4 describes
(stability coefficients in or near published ranges, correct cross-sectional
signatures). Fidelity to published USE4 \emph{exposures} is necessarily reduced
for EYLD and GRO in particular; both audits carry this caveat.

%──────────────────────────────────────────────────────────────────────────────
\section{Register A --- Data Availability}

%──────────────────────────────────
\subsection{A1. EYLD: EPFWD omitted, weights renormalised \hfill \WARN{Material}}

\begin{description}[style=nextline]
  \item[USE4] $\mathrm{EYLD} = 0.75\,\mathrm{EPFWD} + 0.15\,\mathrm{CETOP} +
    0.10\,\mathrm{ETOP}$; EPFWD is the twelve-month forward earnings-to-price
    from analyst-consensus EPS.
  \item[Deviation] No analyst estimates exist in Sharadar; the descriptor that
    carries 75\% of the composite cannot be built.
  \item[Solution] Drop EPFWD and renormalise over the trailing descriptors:
    $\mathrm{EYLD} = 0.60\,\mathrm{CETOP} + 0.40\,\mathrm{ETOP}$ (preserving the
    PDF's 0.15:0.10 ratio between them).
  \item[Validity] The alternative --- proxying EPFWD with trailing EPS --- adds
    no information (it duplicates ETOP) while pretending to forward content;
    renormalisation is the honest choice. The trailing composite still behaves
    as an earnings-yield factor should: ESTU median raw eyld \textbf{0.05891}
    (a realistic US E/P), positive-median on \textbf{100.0\%} of dates,
    quintiles strictly monotone (spread $0.26033$), MoM $\rho$ \textbf{0.9630}.
    What is lost is the forward-looking revision signal, not the value tilt.
  \item[Revisit when] A forward-estimates feed (IBES-class) is procured ---
    then restore $0.75/0.15/0.10$ and rebuild.
\end{description}

%──────────────────────────────────
\subsection{A2. GRO: EGRLF proxied by EGRO \hfill \WARN{Material}}

\begin{description}[style=nextline]
  \item[USE4] $\mathrm{GRO} = 0.70\,\mathrm{EGRLF} + 0.20\,\mathrm{EGRO} +
    0.10\,\mathrm{SGRO}$; EGRLF is the analyst long-term (3--5yr) predicted
    earnings growth.
  \item[Deviation] Same root cause as A1: no analyst estimates.
  \item[Solution] Proxy EGRLF with EGRO (\texttt{GRO\_W\_EGRLF\_PROXY} $=$
    \texttt{GRO\_W\_EGRLF}), giving effective
    $\mathrm{GRO} = 0.90\,\mathrm{EGRO} + 0.10\,\mathrm{SGRO}$.
  \item[Validity] Unlike A1, here a proxy (rather than renormalisation) is
    defensible because EGRLF and EGRO target the same quantity --- earnings
    growth --- at different horizons, and trailing growth is the standard
    instrument for expected growth when estimates are absent. The realised
    factor shows the expected behaviour: quintiles strictly monotone (spread
    $1.25906$), MoM $\rho$ \textbf{0.9863} (annual inputs are slow-moving),
    ESTU median $-0.08120$ within the plausible band. The caveat stands: 70\%
    of the intended exposure is trailing, not forward, so GRO under-captures
    expectations-driven growth repricing.
  \item[Revisit when] Analyst long-term growth estimates are procured --- set
    \texttt{GRO\_W\_EGRLF\_PROXY} $=0$ and rebuild.
\end{description}

%──────────────────────────────────
\subsection{A3. Total vs float-adjusted market cap \hfill \WARN{Material}}

\begin{description}[style=nextline]
  \item[USE4] MSCI USA IMI targets 99\% of \emph{float-adjusted} market cap;
    cap-weighting and $\sqrt{\text{mcap}}$ regression weights follow.
  \item[Deviation] Sharadar provides total market cap only. The gap touches
    ESTU selection and ranks, all CW means, the benchmark, SIZE/NLS exposures,
    and every $\sqrt{\text{mcap}}$ WLS weight.
  \item[Solution] Use total mcap consistently everywhere; never mix.
  \item[Validity] For US large-caps the float fraction is typically 80--100\%,
    so rank and weight distortions are second-order precisely where the cap
    weight is concentrated; closely-held names are penalised toward inclusion
    rather than exclusion (conservative for ESTU membership). The empirical
    fingerprint matches the intended universe: ESTU cap-weighted beta
    \textbf{1.007} ($\approx 1.0$ identity), benchmark correlation to the broad
    CRSP proxy \textbf{0.9333}, mcap coverage of the retain-eligible universe
    mean \textbf{99.3\%} --- the 99\% IMI design target.
  \item[Revisit when] Float-share data is procured --- re-derive
    \texttt{estu\_monthly}, then rebuild downstream (the run order in
    \texttt{run\_pipeline.py} already encodes the cascade).
\end{description}

%──────────────────────────────────
\subsection{A4. Preferred equity unavailable (BP, LEV) \hfill \WARN{Material}}

\begin{description}[style=nextline]
  \item[USE4] BP: $\mathrm{BE} = \text{total equity} - \text{preferred equity}$.
    LEV: $\mathrm{MLEV} = (\mathrm{ME}+\mathrm{PE}+\mathrm{LD})/\mathrm{ME}$,
    $\mathrm{BLEV} = (\mathrm{BE}+\mathrm{PE}+\mathrm{LD})/\mathrm{BE}$.
  \item[Deviation] SF1 has no preferred-equity balance field (only the
    income-statement \texttt{prefdivis}); PE $=0$ throughout.
  \item[Solution] BP uses \texttt{equity} directly as BE. LEV computes
    $\mathrm{BLEV} = (\mathrm{EQUITY}+\mathrm{DEBTNC})/\mathrm{EQUITY}$ --- note
    the PE terms cancel algebraically in BLEV's numerator
    ($\mathrm{BE}+\mathrm{PE} = \mathrm{EQUITY}$), so only BLEV's
    \emph{denominator} and MLEV's numerator are affected.
  \item[Validity] Preferred stock is rare among ESTU members and small where it
    exists, concentrated in financials/utilities; for those names BP's BE is
    slightly overstated and MLEV slightly understated --- both in the direction
    of \emph{damping} the factor, not inverting it. The cross-sections are
    undistorted: BP median \textbf{0.4077} with strictly monotone quintiles;
    LEV bounded ($[-0.564, 2.772]$ post-trim) with monotone quintiles (spread
    $1.75542$) and the heaviest-leverage names (where PE matters most) still
    sorting to Q5. The alternative --- estimating PE from \texttt{prefdivis}
    via an assumed yield --- would inject a noisy model where the field should
    be a balance-sheet fact.
  \item[Revisit when] A preferred-equity balance field becomes available
    (Sharadar schema change or supplemental source) --- restore
    $\mathrm{BE} = \mathrm{EQUITY} - \mathrm{PE}$ in BP and the full MLEV/BLEV
    forms in LEV.
\end{description}

%──────────────────────────────────
\subsection{A5. CETOP proxied by net income + depreciation (EYLD)}

\begin{description}[style=nextline]
  \item[USE4] CETOP is cash earnings to price.
  \item[Deviation] No cash-earnings field in SF1.
  \item[Solution] $\mathrm{CETOP} = (\texttt{netinc} + \texttt{depamor})/\text{mcap}$.
  \item[Validity] Adding back depreciation and amortisation to net income is the
    textbook first-order cash-earnings approximation and uses fields with
    \textbf{0} nulls in the TTM panel. Spot medians are economically sensible
    (median eyld \textbf{0.0332} at 2026-05-29).
  \item[Revisit when] A cash-flow-statement-based cash earnings field is
    preferred --- swap the numerator; the pipeline is otherwise unchanged.
\end{description}

%──────────────────────────────────
\subsection{A6. MSCI IMI membership proxied (ESTU)}

\begin{description}[style=nextline]
  \item[USE4] ESTU $=$ MSCI USA IMI; the construction rules are proprietary.
  \item[Deviation] Membership lists are unavailable; the universe must be
    reconstructed from observables.
  \item[Solution] Domestic common stocks on NYSE/NASDAQ/NYSEMKT, \$300M floor,
    ATVR liquidity screen (add $\geq 20\%$, retain $\geq 10\%$, 63-day window),
    top-3{,}000 cap target with a 3{,}000/3{,}500 hysteresis buffer ---
    mirroring MSCI's published IMI design principles (breadth, investability,
    stability).
  \item[Validity] The proxy reproduces the IMI's measurable properties: size
    \textbf{2{,}187--3{,}045} recent (IMI-like breadth), churn mean
    \textbf{2.3\%} with $95.5\%$ of months $<5\%$ (stability), mcap coverage
    \textbf{99.3\%} of the retain-eligible universe (the 99\% target), segment
    proportions large $11.5\%$ / mid $15.3\%$ / small $73.2\%$, and mega-caps
    included every eligible month. Eight-check battery green
    (\texttt{estu\_audit.tex}).
  \item[Revisit when] Churn exceeds 5\% sustained, or licensed IMI constituents
    become available for a direct overlap study.
\end{description}

%──────────────────────────────────
\subsection{A7. SEP turnover anomalies (LIQ)}

\begin{description}[style=nextline]
  \item[USE4] Share turnover $=$ volume / shares outstanding.
  \item[Deviation] \textbf{20{,}710} SEP rows ($0.057\%$, 2{,}221 tickers) show
    physically impossible turnover $>1$ (vendor share-count artefacts).
  \item[Solution] No row-level surgery; the log transform plus the 3$\sigma$
    trim absorb them.
  \item[Validity] The anomalies are five orders of magnitude too sparse to move
    monthly medians; the unit check confirms the bulk is sound (median daily
    turnover \textbf{0.00461}, inside the 0.001--0.010 plausibility band), and
    the post-trim ESTU distribution is clean (clip rate \textbf{0.95\%}, skew
    $-0.057$). Excising rows would risk deleting genuine high-volume events
    (squeezes) along with artefacts.
  \item[Revisit when] The impossible-row share grows past $\approx 0.5\%$ of SEP
    or clusters in ESTU members.
\end{description}

%──────────────────────────────────
\subsection{A8. TTM dimension absent (EYLD, DYLD)}

\begin{description}[style=nextline]
  \item[USE4] Trailing-twelve-month income items.
  \item[Deviation] The SF1 extract carries dimension ARQ only.
  \item[Solution] Construct TTM as 4-quarter rolling sums over
    \texttt{calendardate} (EYLD: netinc, depamor --- requiring 4 complete
    quarters; DYLD: split-adjusted DPS with \texttt{MIN\_QUARTERS}=1 so
    zero-payers survive).
  \item[Validity] A 4-quarter sum over fiscal quarters \emph{is} TTM by
    definition; requiring 4 complete quarters (EYLD) avoids annualising
    partial years. Evidence: \textbf{567{,}298} clean TTM rows (EYLD),
    \textbf{608{,}164} (DYLD), zero negative TTM DPS, and downstream medians in
    economic ranges.
  \item[Revisit when] Sharadar's TTM dimension is licensed --- a direct
    cross-check of the constructed sums would then be run before switching.
\end{description}


%──────────────────────────────────
\subsection{A9. GICS unavailable --- engineered 55-factor scheme \hfill \WARN{Material}}

\begin{description}[style=nextline]
  \item[USE4] 60 industry factors built on GICS --- itself a hybrid recombination
    of GICS levels selected on economic intuition, statistical significance,
    explanatory power, and minimum size.
  \item[Deviation] GICS is licensed and absent from Sharadar.
  \item[Solution] Sharadar's \texttt{industry} field (154 observed atoms, 100\%
    ESTU coverage) as the atomic layer, aggregated by the hand-engineered,
    version-controlled \path{industry_scheme.csv} into \textbf{55 factors}
    targeting the published 60-list under the published \emph{criteria}.
  \item[Validity] USE4 itself states the factor set is engineered for stability
    and explanatory power, ``not taxonomy purity'' --- replicating the 2011
    labels verbatim in 2026 would violate the spec's own criteria after fifteen
    years of sector tectonics. The empirical fingerprint is sound: all 55
    factors populated at every one of 330 dates, largest factor never exceeds
    \textbf{14.1\%} of ESTU cap, floor violations \textbf{0}, and all six
    mega-cap spot checks land. The scheme is a reviewable 154-row table, not
    code.
  \item[Revisit when] GICS (or a maintained crosswalk) is licensed --- run an
    overlap study before switching.
\end{description}

%──────────────────────────────────
\subsection{A10. No segment data --- single 0/1 membership \hfill \WARN{Material}}

\begin{description}[style=nextline]
  \item[USE4] Diversified firms receive fractional industry exposures from
    business-segment assets and sales (blended 75\%/25\%, top five segments);
    $\sim$63\% of the estimation universe by cap weight was multi-exposed.
  \item[Deviation] SF1 carries no business-segment reporting; the fractional
    machinery is unbuildable from public Sharadar.
  \item[Solution] Single 0/1 membership for every stock; the raw classification
    atom is retained as an audit column.
  \item[Validity] Single membership is the standard academic treatment and
    leaves the CSR identification (cap-weighted zero-sum industry constraint)
    untouched --- exposures still sum to 1 per stock by construction. The bias
    is named and concentrated: conglomerates, diversified financials, and
    integrated energy are forced into one bucket each, and the affected factors
    (Industrial Conglomerates, Consumer Finance \& DFS, Oil Gas \& Consumable
    Fuels) are exactly where the audit flags interpretation caveats. Hand-curating
    fractional exposures for the top conglomerates was considered and rejected:
    a subjective, unreproducible input is worse than a documented uniform rule.
  \item[Revisit when] Segment-level fundamentals are procured.
\end{description}

%──────────────────────────────────
\subsection{A11. Classification is not point-in-time}

\begin{description}[style=nextline]
  \item[USE4] Classifications maintained point-in-time by MSCI's internal
    process.
  \item[Deviation] Sharadar TICKERS is a current snapshot --- today's label
    applied to 1999 holdings carries mild look-ahead for reclassified or
    pivoted companies.
  \item[Solution] Static classification, documented; the build asserts the map
    is static (one label per permaticker, ever), so any future PIT layer is a
    drop-in replacement.
  \item[Validity] Industry membership drifts at a far lower frequency than any
    factor signal --- reclassification is a multi-year event, and the affected
    names are a thin tail of the panel. This is the standard compromise in
    academic work built on snapshot classifications. A genuinely point-in-time
    path exists at zero licence cost (EDGAR 10-K header SIC codes are
    filing-dated), so the deviation is a deferral, not a dead end.
  \item[Revisit when] Classification drift is shown to matter (e.g., a
    dot-com-era membership study) --- build the EDGAR PIT layer.
\end{description}

%──────────────────────────────────────────────────────────────────────────────
\section{Register B --- Methodology Judgment Calls}

%──────────────────────────────────
\subsection{B1. PIT key and staleness window (all fundamentals)}

\begin{description}[style=nextline]
  \item[USE4] ``Most recent'' / ``last reported'' values; no operational PIT
    rule given.
  \item[Deviation] Two free choices: which date stamps a filing, and how old a
    filing may be.
  \item[Solution] \texttt{datekey} (the filing release date) as the PIT key ---
    never \texttt{calendardate} --- joined backward
    (\texttt{datekey} $\leq$ \texttt{signal\_date}); staleness cutoff
    \texttt{LOOKBACK\_DAYS} $=548$ ($\approx 18$ months) family-wide.
  \item[Validity] \texttt{calendardate} would leak 30--90 days of look-ahead
    per filing (results are filed weeks after the period closes) --- this is
    the single most common factor-construction bug and the pipeline forecloses
    it structurally. The 548-day window tolerates late filers (annual report
    $+$ up to 6 months delay) while bounding zombie data; coverage validates at
    \textbf{4{,}097--7{,}261} names post-2005 across the fundamental factors.
  \item[Revisit when] Never for the key; the window if coverage floors are
    threatened (see C6).
\end{description}

%──────────────────────────────────
\subsection{B2. DYLD corrupt-DPS guard and split adjustment}

\begin{description}[style=nextline]
  \item[USE4] Trailing dividends per share over price; no data-quality protocol.
  \item[Deviation] Raw vendor DPS contains split-era artefacts implying absurd
    yields.
  \item[Solution] Two-layer \texttt{DPS\_Q\_YIELD\_MAX} $=0.50$ guard ---
    quarterly DPS implying a single-quarter yield $>50\%$ is nulled at source
    (\textbf{1{,}292} filings); the same threshold backstops the TTM yield
    (\textbf{8{,}700} nulled). DPS re-based across \textbf{217{,}223}
    split-affected quarters via cumulative split factors.
  \item[Validity] A 50\% \emph{quarterly} yield is not a dividend policy that
    exists; it is a share-count mismatch. Nulling (rather than capping) keeps
    fabricated values out of the TTM sum entirely. Post-guard the factor is
    clean: zero negative dyld, max post-trim \textbf{0.2020} (well under the
    ceiling), median \textbf{0.00285} --- a realistic US yield --- and MoM
    $\rho$ \textbf{0.9927}.
  \item[Revisit when] The flagged-filing count grows materially after a vendor
    refresh (currently $0.2\%$ of ARQ rows).
\end{description}

%──────────────────────────────────
\subsection{B3. Negative book equity retained (BP)}

\begin{description}[style=nextline]
  \item[USE4] Silent on negative BE.
  \item[Deviation] \textbf{49{,}043} ARQ rows carry negative equity.
  \item[Solution] Retain as negative BTOP; null only when \texttt{equity} is
    null.
  \item[Validity] Negative book equity is real information --- buy-back-heavy
    and distressed/levered firms --- and the value cross-section should rank
    them below low-BTOP growth names, which negative values do automatically.
    Marking them missing would silently delete a coherent economic cohort.
    Validation asserts the property explicitly (check 2b: negative btop present
    post-2000, \PASS{PASS}).
  \item[Revisit when] Never; this is the economically correct treatment.
\end{description}

%──────────────────────────────────
\subsection{B4. Zero-payers are zeros, not missing (DYLD)}

\begin{description}[style=nextline]
  \item[USE4] Silent.
  \item[Deviation] $\approx 63\%$ of covered names pay no dividend.
  \item[Solution] A confirmed payer of nothing gets \texttt{dyld} $=0.0$
    (with \texttt{MIN\_QUARTERS}=1); only stocks with no filings at all are
    missing.
  \item[Validity] ``Pays no dividend'' is a statement about the company;
    ``no data'' is a statement about the database --- conflating them would
    shrink coverage by nearly two thirds and turn DYLD into a payer-only
    factor. The quintile structure confirms the design: Q1/Q2 means exactly
    \textbf{0.00000} (the zero-payer mass), Q3--Q5 strictly increasing to
    \textbf{0.05579}.
  \item[Revisit when] Never; definitional.
\end{description}

%──────────────────────────────────
\subsection{B5. EGRO/SGRO pre-clip at $\pm 5$ (GRO)}

\begin{description}[style=nextline]
  \item[USE4] Normalised slope (slope / mean EPS); no guard specified.
  \item[Deviation] The normalisation explodes when mean EPS is tiny but
    positive.
  \item[Solution] Clip EGRO and SGRO to $[-5,+5]$ before compositing
    (\textbf{41{,}815} EGRO rows, \textbf{1{,}673} SGRO rows on this run).
  \item[Validity] $\pm 500\%$ annualised growth is beyond any economically
    meaningful signal; values outside it are denominators, not growth. Without
    the pre-clip these rows would contaminate the 3$\sigma$ bounds themselves
    (the trim would be set by artefacts). With it, the trim operates on a sane
    distribution and the quintile spread ($1.25906$) stays monotone.
  \item[Revisit when] Never in principle; the constant if clipping ever
    exceeds a few percent of rows.
\end{description}

%──────────────────────────────────
\subsection{B6. Null debt treated as zero (LEV)}

\begin{description}[style=nextline]
  \item[USE4] Silent on null fields.
  \item[Deviation] \texttt{debtnc} is null on \textbf{20.3\%} of ARQ rows.
  \item[Solution] \texttt{DEBT\_FLOOR}: null (and any negative) debt $\to 0$ in
    the MLEV/DTOA/BLEV numerators.
  \item[Validity] For balance-sheet debt, absence of the line item means the
    firm carries none --- null is ``zero'', not ``unknown'' (the same firms
    report assets and equity on the same filing). Treating it as missing would
    delete one fifth of the universe, biased toward exactly the unlevered
    firms LEV must rank lowest. Sanity holds: min MLEV/BLEV $=1.0000$ and
    DTOA $>1$ on only $1.2\%$ of rows.
  \item[Revisit when] Never; accounting semantics.
\end{description}

%──────────────────────────────────
\subsection{B7. Standardise-then-combine (LEV)}

\begin{description}[style=nextline]
  \item[USE4] $0.75\,\mathrm{MLEV} + 0.15\,\mathrm{DTOA} + 0.10\,\mathrm{BLEV}$
    --- units unstated.
  \item[Deviation] The raw sub-descriptors live on incompatible scales
    (MLEV/BLEV $\in [1, \approx 20]$; DTOA $\in [0, 1]$).
  \item[Solution] Standardise each sub-descriptor (CW mean 0 / EW std 1 on
    ESTU) before applying the weights.
  \item[Validity] A raw-scale weighted sum would hand MLEV far more than its
    intended 75\% share --- the weights are only meaningful on a common scale,
    and USE4's own standardisation convention supplies exactly that scale.
    Verified on-run: each sub-descriptor's CW mean is machine-zero
    ($5.96/1.78/1.66 \times 10^{-16}$) before compositing.
  \item[Revisit when] Never; required for the published weights to mean what
    they say.
\end{description}

%──────────────────────────────────
\subsection{B8. SIZE trim: $4\sigma$ around the EW mean}

\begin{description}[style=nextline]
  \item[USE4] Generic outlier rule: trim to three standard deviations from the
    mean.
  \item[Deviation] Applying the family-standard CW-centred $3\sigma$ band to
    \texttt{lncap} would clip mega-caps wholesale: the CW mean of log-cap sits
    $\approx 2\sigma$ above the EW median by construction.
  \item[Solution] For SIZE only: $4\sigma$ band centred on the \emph{EW} mean.
  \item[Validity] The generic rule presumes the trim centre is representative
    of the cross-section; for log market cap the CW centre is definitionally
    not (it is the object the factor measures). The EW-centred $4\sigma$ band
    preserves the genuine right tail --- mega-caps are correct data, not
    outliers --- while still catching artefacts: clip rate \textbf{0.05\%}
    (404 of 822{,}918 ESTU rows), and the standardisation identity still holds
    exactly (CW mean of \texttt{size\_score} $3.37\times10^{-15}$).
  \item[Revisit when] Never in direction; the width if data errors ever appear
    inside $4\sigma$.
\end{description}

%──────────────────────────────────
\subsection{B9. Outlier method for the non-linear factors (NLB trim, NLS winsor)}

\begin{description}[style=nextline]
  \item[USE4] ``the factor is winsorized'' --- percentiles unspecified, for
    both NLB and NLS.
  \item[Deviation] One free choice per factor; the two factors resolve it
    differently.
  \item[Solution] NLS: percentile winsorisation at 1\%/99\% (ESTU quantiles,
    applied to all). NLB: the family-standard $3\sigma$ trim.
  \item[Validity] The choice follows each input's shape. \texttt{nls\_raw} is a
    skewed cubic for which $\sigma$-based bounds produce near-zero clipping
    (the std is inflated by the very tail being treated) --- percentile
    winsorisation clips a designed $2.00\%$ exactly. \texttt{nlb\_raw}'s
    distribution is symmetric enough for the $\sigma$ rule to bind
    (\textbf{2.42\%} ESTU clip, flagged \WARN{PARTIAL MATCH} in its audit) and
    keeping the family trim preserves comparability. Both satisfy the final
    standardisation identity, and both factors land in their USE4-published
    stability ranges (NLS $\rho$ \textbf{0.9882} vs published 0.98--0.99; NLB
    \textbf{0.8723}).
  \item[Revisit when] NLB's clip rate drifts past $\approx 3\%$ --- switch NLB
    to 1\%/99\% winsorisation for symmetry with NLS.
\end{description}

%──────────────────────────────────
\subsection{B10. Orthogonalisation choices (NLB, NLS, RESVOL)}

\begin{description}[style=nextline]
  \item[USE4] ``orthogonalized \ldots on a regression-weighted basis''
    (weights $\propto \sqrt{\text{mcap}}$); pairwise vs joint unstated;
    RESVOL ``orthogonalized to Beta''.
  \item[Deviation] Operational details are free: regression form, fit
    universe, application universe, and (for RESVOL) the regressor and order
    of operations.
  \item[Solution] Pairwise $\sqrt{\text{mcap}}$-weighted WLS with intercept,
    fit on ESTU, residual applied to all stocks (NLB on beta, NLS on
    \texttt{log\_mcap}). RESVOL: \emph{no-intercept} WLS of the post-trim
    composite on \texttt{beta\_score}, per date.
  \item[Validity] Pairwise is Appendix~A's literal reading; joint
    orthogonalisation belongs to the full cross-sectional regression where it
    will happen anyway. The defining property is enforced to machine
    precision on-run: weighted inner products $3.40\times10^{-15}$ (NLB),
    $8.00\times10^{-13}$ (NLS), $1.2041\times10^{-17}$ (RESVOL) --- all
    against a $10^{-10}$ gate. RESVOL's no-intercept form is correct because
    \texttt{beta\_score} is already CW-mean-zero, so an intercept would simply
    re-absorb the level the final standardisation removes anyway. Downstream
    signatures confirm intent: NLS market correlation $-0.1331$
    ($|\cdot|<0.20$, size-neutral barbell); RESVOL retains the low-vol anomaly
    ($0.3532$).
  \item[Revisit when] The multi-factor CSR is assembled --- re-express
    orthogonalisation jointly there and keep these pairwise versions as
    standalone-factor definitions.
\end{description}

%──────────────────────────────────
\subsection{B11. Factor-on-factor inputs from deliverables (NLB, NLS, RESVOL)}

\begin{description}[style=nextline]
  \item[USE4] NLB/NLS cube ``the standardized exposure''; RESVOL's HSIGMA uses
    the Beta regression residuals.
  \item[Deviation] The upstream exposure could be recomputed locally (as early
    versions did) or read from the upstream deliverable.
  \item[Solution] Read the deliverable: NLB consumes
    \path{beta_use4.parquet}; NLS consumes \path{size_use4.parquet}
    (\texttt{log\_mcap} $:=$ post-trim \texttt{lncap}, \texttt{size\_z} $:=$
    \texttt{size\_score}); RESVOL consumes the Beta artifact for HSIGMA and
    the orthogonalisation.
  \item[Validity] The factor being orthogonalised \emph{against} must be the
    factor that enters the risk model, or the collinearity NLB/NLS exist to
    remove is only partially removed. Local recomputation (pre-trim,
    pre-standardisation drift) silently breaks that identity; reading the
    deliverable makes it exact by construction and adds a free consistency
    property: one Beta, one Size, everywhere. Cost is an explicit build order
    (\texttt{beta} $\to$ \texttt{nlb}, \texttt{size} $\to$ \texttt{nls}),
    enforced by \texttt{run\_pipeline.py} and loud loader asserts.
  \item[Revisit when] Never; this is strictly more correct than recomputation.
\end{description}

%──────────────────────────────────
\subsection{B12. Estimator constants not in the PDF (BETA, MOM, RESVOL)}

\begin{description}[style=nextline]
  \item[USE4] Windows and half-lives are published (Beta 252d/63d HL; MOM 504d
    window, 21d lag, 126d HL; DASTD 252d/42d, HSIGMA 252d/63d, CMRA 12
    months); minimum-observation rules are not.
  \item[Deviation] \texttt{MIN\_OBS} must be invented.
  \item[Solution] Beta 63 ($\approx$ one quarter of dailies); MOM and RESVOL
    252 (a full year --- their estimators average over the whole window).
  \item[Validity] The floor trades coverage against estimator variance: a
    regression slope stabilises with one quarter of points, while a weighted
    mean/range over two years is meaningless on a part-year sample. The
    outcomes sit in published ranges --- Beta MoM $\rho$ \textbf{0.9458}
    (USE4 reports 0.96--0.97), MOM \textbf{0.8702} (published 0.89--0.92),
    with coverage floors cleared at \textbf{4{,}854} / \textbf{4{,}701} /
    \textbf{4{,}021}.
  \item[Revisit when] Coverage floors approach binding (see C6) ---
    \texttt{MIN\_OBS} is the first relief valve.
\end{description}

%──────────────────────────────────
\subsection{B13. Benchmark and risk-free conventions (daily artifact)}

\begin{description}[style=nextline]
  \item[USE4] Excess returns over the estimation-universe cap-weighted return;
    RF series unspecified.
  \item[Deviation] Both series must be sourced.
  \item[Solution] Ken French daily RF; benchmark $=$ ESTU cap-weighted excess
    return with \texttt{mcap\_lag} weights (weight known before the return
    accrues), built once in \texttt{daily\_build.ipynb}.
  \item[Validity] Over any 252-day window the RF is nearly constant, so slope
    estimators are insensitive to the exact series --- fidelity, not
    sensitivity, motivates including it. Lagged weights remove the
    contemporaneous-weight bias. The benchmark validates as a market:
    correlation \textbf{0.9333} to the broad CRSP proxy, vol \textbf{19.5\%}
    full-sample with GFC \textbf{55.5\%} / COVID \textbf{69.0\%} crisis peaks,
    \textbf{0} missing dates.
  \item[Revisit when] Only if validating against USE4's published factor
    returns requires their exact universe.
\end{description}


%──────────────────────────────────
\subsection{B14. Industry scheme engineering: floor, merges, splits}

\begin{description}[style=nextline]
  \item[USE4] ``Minimum size'' is a selection criterion; values unpublished.
    The 60-list embeds 2011 merge/split decisions.
  \item[Deviation] The floor values, the four thin-factor exceptions, and every
    merge/split versus the 60-list are our calls.
  \item[Solution] Floor: median $\geq 15$ and min $\geq 8$ ESTU members over the
    full sample. Exceptions (4, declared): Internet \& Catalog Retail (median 6
    --- but \$3.0T of cap: Amazon), Managed Health Care (USE4-explicit),
    Airlines, Industrial Conglomerates. Merges: Drilling$\to$Equipment \&
    Services; Paper$+$Forest$\to$Construction Materials; both metals factors
    $\to$ Metals \& Mining; durables/apparel/leisure and retail consolidated;
    telecom $2\to1$ (no wireless split in the atoms). Splits: Diversified
    Financials and Software each split in two.
  \item[Validity] The floor operationalises USE4's own stability criterion for a
    $\sim$2{,}500-name universe; the exceptions are justified where member
    counts mislead (cap weight is the relevant size for a cap-weighted
    regression). The splits address 2026 realities USE4 could not see ---
    Software and Diversified Financials were 107--151-median-member monsters as
    single factors, larger than entire sectors elsewhere. Validation: \textbf{0}
    floor violations outside the declared exceptions; all 55 factors present at
    every date.
  \item[Revisit when] The CSR runs --- explanatory-power and VIF screening on
    actual factor returns replaces member counts as the selection statistic
    (USE4's remaining two criteria, deferred honestly).
\end{description}

%──────────────────────────────────────────────────────────────────────────────
\section{Register C --- Engineering and Validation Conventions}

%──────────────────────────────────
\subsection{C1. Shared daily artifact and memory gate}

\begin{description}[style=nextline]
  \item[USE4] Not addressed (implementation concern).
  \item[Deviation] Five time-series builds previously re-downloaded RF and
    re-loaded/sorted the $\approx$39M-row price panel independently; one OOM
    kernel death traced to concurrent commit exhaustion.
  \item[Solution] \texttt{daily\_build.ipynb} builds the panel once
    (\textbf{38{,}997{,}999} rows, 377.9 MB zstd); factor builds load it via
    \texttt{factor\_lib.load\_daily\_artifacts} behind a commit-aware
    FAST/SEQUENTIAL gate; a FAST$\equiv$SEQUENTIAL equivalence check
    ($<10^{-12}$) guards the two load paths.
  \item[Validity] Numerically a no-op (the artifact is byte-identical input to
    what each build derived); operationally it removes the network dependency,
    runs the equivalence check once instead of five times, and eliminates the
    crash mode by construction. The run scorecard (\S2) is the evidence: all
    five consumers green from one artifact.
  \item[Revisit when] You productionise (the daily-panel loader is the first
    thing to move into a real package).
\end{description}

%──────────────────────────────────
\subsection{C2. Completed-months coverage convention}

\begin{description}[style=nextline]
  \item[USE4] Not addressed (validation concern).
  \item[Deviation] The final signal date is the in-progress month: freshest
    prices/fundamentals lag the vendor refresh, so its coverage dips below any
    completed month and intermittently failed fixed floors.
  \item[Solution] Check 3 in every factor build evaluates the coverage floor on
    completed months only (the final date excluded), labelled as such.
  \item[Validity] The in-progress date measures vendor latency, not pipeline
    coverage; a floor that fails on Tuesday and passes on month-end is testing
    the calendar. The first incident (RESVOL: 3{,}983 on the partial date vs
    \textbf{4{,}021} completed-month minimum) showed exactly this signature ---
    every completed month was above floor.
  \item[Revisit when] Never for the convention; the floors per C6.
\end{description}

%──────────────────────────────────
\subsection{C3. Coverage rule: ESTU moments, universe-wide scores}

\begin{description}[style=nextline]
  \item[USE4] Standardisation moments from the estimation universe; exposures
    needed for the full coverage universe.
  \item[Deviation] None in substance --- registered because it explains
    recurring ``low'' all-universe statistics.
  \item[Solution] CW mean / EW std computed on ESTU only, applied to all
    stocks; all reports show the in-band fraction for \emph{both} populations.
  \item[Validity] This is USE4's own coverage design; its visible consequence
    is that non-ESTU micro-caps legitimately standardise outside $[\pm 3]$
    (e.g.\ LIQ: all-universe \textbf{84.3\%} vs ESTU \textbf{98.9\%}; SIZE:
    \textbf{44.4\%} vs \textbf{80.4\%}). Reporting both prevents the
    structural gap from masquerading as a defect.
  \item[Revisit when] Never.
\end{description}

%──────────────────────────────────
\subsection{C4. Family trim machinery and accepted out-of-band clip rates}

\begin{description}[style=nextline]
  \item[USE4] Trim to three standard deviations from the mean.
  \item[Deviation] The PDF's mean/std are unspecified; and three factors run
    hot against the working 0.5--2\% clip-rate band.
  \item[Solution] Family standard: CW mean $\pm\,3\cdot$EW std per date, ESTU
    bounds, clip applied to all stocks. Out-of-band rates are kept and flagged
    \WARN{PARTIAL MATCH} in the audits: GRO \textbf{2.74\%}, LEV
    \textbf{2.71\%}, NLB \textbf{2.42\%}.
  \item[Validity] The CW-mean centre matches the standardisation that follows
    (the same $\mu$ the score subtracts), making trim and score consistent.
    The hot clip rates are distributional facts --- fat-tailed growth slopes,
    leverage's heavy right tail, cubic tails --- not parameter errors;
    widening the band to force 2\% would simply admit the tails the trim
    exists to control. Honest flagging beats silent retuning.
  \item[Revisit when] Any rate passes $\approx 3\%$ --- then reconsider the
    method (percentile winsorisation) per factor, as B9 anticipates for NLB.
\end{description}

%──────────────────────────────────
\subsection{C5. Dual-class deduplication}

\begin{description}[style=nextline]
  \item[USE4] MSCI treats share classes by float; not reproducible without
    float data.
  \item[Deviation] Sharadar lists each class with its own row per day under
    one permaticker in the price panel.
  \item[Solution] Keep the highest-mcap row per \texttt{(permaticker, date)}
    at load time, uniformly in \texttt{estu} and \texttt{daily}.
  \item[Validity] Deterministic, PIT-safe, and conservative: the primary
    economic listing represents the issuer, and double counting (both classes
    as separate cap weight under one issuer id) is avoided. Mega-cap spot
    months (GOOGL 261) confirm continuity across class structures.
  \item[Revisit when] Float data arrives (A3) --- per-class treatment becomes
    possible.
\end{description}

%──────────────────────────────────
\subsection{C6. Coverage floors vs the shrinking listed universe}

\begin{description}[style=nextline]
  \item[USE4] Not addressed.
  \item[Deviation] Coverage floors (4{,}000; LEV 3{,}800) were set against a
    larger US listed universe; post-2005 minima now run 4{,}021--7{,}261 and
    drift down with delistings.
  \item[Solution] Keep the floors, monitor the margins; RESVOL's audit states
    its \textbf{0.5\%} margin explicitly.
  \item[Validity] The floors are tripwires, not requirements of USE4; their
    value is alerting on \emph{discontinuous} coverage loss (a broken join)
    rather than secular drift. Keeping them tight maximises that sensitivity;
    the documented margins distinguish drift from breakage when they fire.
  \item[Revisit when] Any completed-month minimum breaches its floor ---
    lower the floor deliberately (with the LEV-style documented reason) rather
    than reflexively.
\end{description}


%──────────────────────────────────
\subsection{C7. Fail-loudly guard on unmapped classification atoms}

\begin{description}[style=nextline]
  \item[USE4] Not addressed (implementation concern).
  \item[Deviation] Vendor taxonomies grow; a silent default bucket would let
    unreviewed atoms leak into the exposure matrix.
  \item[Solution] \texttt{industries\_build.ipynb} asserts every observed atom
    is present in \path{industry_scheme.csv} and fails the build otherwise ---
    taxonomy changes are reviewed, never defaulted.
  \item[Validity] The guard fired on its first execution: three atoms
    (Banks--Global, Drug Manufacturers--Major, Silver) exist in the full ticker
    universe but never among ESTU members, so the ESTU-derived design worksheet
    missed them. Each was reviewed and assigned in one line of the scheme. A
    silent fallback would have buried exactly this class of gap.
  \item[Revisit when] Never; the assert is the mechanism.
\end{description}

%──────────────────────────────────────────────────────────────────────────────
\section{Consolidated Revisit Triggers}

\begin{longtable}{@{}p{1.2cm}p{5.6cm}p{6.8cm}@{}}
\caption{Open triggers, by register entry}\\
\toprule
\textbf{Entry} & \textbf{Watching} & \textbf{Action when fired} \\
\midrule
\endhead
A1, A2 & Analyst-estimates feed procured & Restore EPFWD ($0.75/0.15/0.10$) in EYLD; set EGRLF proxy weight to 0 in GRO; rebuild both \\[4pt]
A3, C5 & Float-share data procured & Re-derive ESTU on float-adjusted mcap; cascade rebuild; revisit per-class handling \\[4pt]
A4 & Preferred-equity field available & Restore $\mathrm{BE}=\mathrm{EQUITY}-\mathrm{PE}$ (BP) and full MLEV/BLEV (LEV) \\[4pt]
A6 & ESTU churn $>5\%$ sustained & Re-tune buffer/ATVR; overlap study if IMI constituents licensed \\[4pt]
A7 & SEP impossible-turnover share $>0.5\%$ & Row-level cleaning in LIQ Stage 1 \\[4pt]
B9, C4 & Any ESTU clip rate $>3\%$ & Switch that factor to 1\%/99\% winsorisation (NLB first candidate) \\[4pt]
B10 & Multi-factor CSR assembled & Joint orthogonalisation in the CSR; pairwise versions remain the standalone definitions \\[4pt]
B12, C6 & Completed-month coverage minimum breaches its floor & Lower floor with documented reason; \texttt{MIN\_OBS} relief if estimator-side \\[4pt]
C1 & Productionising & Move shared loaders/validation helpers into a real package \\
A9 & GICS or crosswalk licensed & Overlap study; re-map the scheme \\[4pt]
A10 & Segment data procured & Build fractional exposures (75/25 assets/sales) \\[4pt]
A11 & Classification drift shown to matter & EDGAR 10-K SIC point-in-time layer \\[4pt]
B14 & CSR running & Re-screen the industry scheme on explanatory power / VIF \\
\bottomrule
\end{longtable}

%──────────────────────────────────────────────────────────────────────────────
\section{Companion Documents}

Per-component evidence backing every number above: \texttt{estu\_audit.tex},
\texttt{daily\_summary.tex}, \texttt{<factor>\_factor\_audit.tex} in each style
directory, and \texttt{industries\_factor\_audit.tex} (with the scheme itself in
\path{03_industry_factors/industry_scheme.csv}) --- all from the same
2026-06-12 reference run audited here.
Construction rationale at spec level: \texttt{<factor>\_risk\_factor\_guide.ipynb}
(NOT-IN-PDF master lists), \texttt{estu/estu\_construction\_guide.ipynb},
\newline\texttt{daily/daily\_construction\_guide.ipynb}.

\end{document}